\documentclass[14pt]{extarticle}
\usepackage[
 a4paper, mag=1000, nohead, includefoot,
 left=3cm, right=1cm, top=2cm, bottom=2cm, footskip=1cm
]{geometry}

% Times New Roman
%\usepackage{times}
\usepackage{mathptmx}
\renewcommand{\rmdefault}{ftm}

%\linespread{1.3}    %полуторный отступ
\usepackage[onehalfspacing]{setspace} %``умное'' расстояние между строк - установить
    %1.5 интервала от нормального, эквивалентно
    %\renewcommand{\baselinestretch}{1.24}

% Пакети перекодування
\usepackage[T2A]{fontenc}
\usepackage[utf8]{inputenc}
\usepackage[russian,ukrainian]{babel}

%\usepackage{amssymb,amsfonts,amsmath,mathtext,cite,enumerate,float} %подключаем нужные пакеты расширений


\usepackage{amssymb,amsmath,amsfonts,latexsym,mathtext} %расширенные наборы математических символов

\usepackage{cite}  %``умные'' библиографические ссылки

\usepackage{enumerate}  %создание и автоматическая нумерация списков

\righthyphenmin=2 % Минимальное число символов при переносе - 2.
% Разного рода графика
%\usepackage{pstricks} % pstricks-графика
%\usepackage{pgf}

% Грамотное PDF vs DVI
\usepackage{ifpdf}

% Подключение графики
\ifpdf %if using pdfLaTeX in PDF mode
  \usepackage[pdftex]{graphicx}
  \DeclareGraphicsExtensions{.pdf,.png,.jpg,.jpeg,.mps}
  \usepackage{pgf}
  \usepackage{tikz}
\else %if using LaTeX or pdfLaTeX in DVI mode
  \usepackage{graphicx}
  \DeclareGraphicsExtensions{.eps,.bmp}
  \DeclareGraphicsRule{.emf}{bmp}{}{}% declare EMF filename extension
  \DeclareGraphicsRule{.png}{bmp}{}{}% declare PNG filename extension
  \usepackage{pgf}
  \usepackage{tikz}
  \usepackage{pstricks}
\fi
\usepackage{epic,bez123}
%\usepackage{floatflt}% package for floatingfigure environment
\usepackage{wrapfig}% package for wrapfigure environment

\graphicspath{{images/}}

% Настройка ссылок в генерируемом pdf  файле.
\usepackage{hyperref}
\hypersetup{
            colorlinks         = false,
            bookmarksopen      = true,
            bookmarksnumbered  = true,
            bookmarksopenlevel = 2,
            pdfstartview       = {FitH},
            pdfborder          = {0 0 1},
            pdfauthor={Bondarchuk Alexey},
            %pdftitle={Купсова Робота},
            %pdfcreator={pdftex},
            %pdfsubject={Курсова робота}
}


\usepackage{indentfirst} % оступ в начале каждого абзаца

% Максимальный уровень вложенности секций, помещаемых в Содержание (по умолчанию вносятся \part, \chapter и \section).
% tocdepth
\setcounter{tocdepth}{3}    

% Максимальный уровень вложенности нумеруемых секций (по умолчанию нумеруются \part,\chapter, \section и \subsection).
% secnumdepth
\setcounter{secnumdepth}{3}


% Добавление в оглалвелние ненумерованых разделов
%\section*{Введение}
%\addcontentsline{toc}{section}{\tocsecindent{Введение}}
\newcommand{\tocsecindent}{\hspace{0mm}}

% Подчинение нумерации формул, рисунков, таблиц
\numberwithin{equation}{section}
\numberwithin{figure}{section}
\numberwithin{table}{section}


% Отбивки возле формул
%\abovedisplayskip=1.3\abovedisplayskip

%\belowdisplayskip=1.3\belowdisplayskip

% переопределение некоторых сущьностей
\makeatletter

%\renewcommand\appendix{\par
%\setcounter{section}{0}%
%      \setcounter{subsection}{0}%
%            \gdef\thesection{\@Alph\c@section}}

%\renewcommand{\@dotsep}{4.5}

% Вид section в теле документа
%\renewcommand{\section}{\@startsection{section}{1}{\parindent}%
%{-3.5ex plus -1ex minus -.2ex}{2.3ex plus.2ex}%
%{\normalfont}}

% Вид subsection в теле документа
%\renewcommand{\subsection}{\@startsection{subsection}{2}{\parindent}%
%{-3.5ex plus -1ex minus -.2ex}{2.3ex plus.2ex}%
%{\normalfont}}

% Вид subsubsection в теле документа
%\renewcommand{\subsubsection}{\@startsection{subsubsection}{3}{\parindent}%
%{-3.5ex plus -1ex minus -.2ex}{2.3ex plus.2ex}%
%{\normalfont}}

\renewcommand\section{\@startsection {section}{1}{\parindent}%
                                   {5.5ex \@plus -1ex \@minus -.2ex}%
                                   {2.3ex \@plus.2ex}%
%                                   {\normalfont\bfseries}}
                                    {\normalfont}}
\renewcommand\subsection{\@startsection{subsection}{2}{\parindent}%
                                     {5.25ex\@plus -1ex \@minus -.2ex}%
                                     {1.5ex \@plus .2ex}%
%                                     {\normalfont\bfseries}}
                                        {\normalfont}}
\renewcommand\subsubsection{\@startsection{subsubsection}{3}{\parindent}%
                                     {5.25ex\@plus -1ex \@minus -.2ex}%
                                     {1.5ex \@plus .2ex}%
%                                     {\normalfont\normalsize\bfseries}}
                                    {\normalfont}}



% Вид section в заголовке

\renewcommand*\l@section[2]{%
  \ifnum \c@tocdepth >\z@
    \addpenalty\@secpenalty
    \addvspace{0.0em \@plus\p@}%
    \setlength\@tempdima{0.8em}%
    \begingroup
      \parindent \z@ \rightskip \@pnumwidth
      \parfillskip -\@pnumwidth
      \leavevmode 
      \advance\leftskip\@tempdima
      \hskip -\leftskip
      #1\nobreak \leaders\hbox to .7em{\hss.\hss}\hfil\nobreak\hb@xt@\@pnumwidth{\hss #2}\par
    \endgroup
  \fi}

\renewcommand*\l@subsection{\@dottedtocline{2}{0.8em}{1.6em}}
\renewcommand*\l@subsubsection{\@dottedtocline{3}{2.4em}{2.4em}}

% Рисунки и таблицы

% переопределение названия, разделительного знака и отступов
\usepackage{caption} 
\DeclareCaptionLabelSeparator{tire}{~--- } 
\captionsetup{justification=centering,labelsep=tire,aboveskip=00pt,belowskip=14pt}

% вариант по ``суворому челябинскому программисту''

% размер отступов
%\setlength\abovecaptionskip{2\p@}
%\setlength\belowcaptionskip{1\p@}

%\def\capfigure{figure}
%\def\captable{table}
%\long\def\@makecaption#1#2{%
%  \vskip\abovecaptionskip
%    \ifx\@captype\capfigure
%        \centering #1~---~#2 \par
%    \else
%        #1~---~#2 \par
%    \fi
%    \vskip\belowcaptionskip
%}

% вариант по ``Львовскому''
%\long\def\@makecaption#1#2{%
%  \vskip\abovecaptionskip
%  \sbox\@tempboxa{#1~--- #2}%
%  \ifdim \wd\@tempboxa >\hsize
%    #1~--- #2\par
%  \else
%    \global \@minipagefalse
%    \hb@xt@\hsize{\hfil\box\@tempboxa\hfil}%
%  \fi
%\vskip\belowcaptionskip}

% названия объектов
\addto\captionsukrainian{\renewcommand\figurename{Рисунок}}
\addto\captionsukrainian{\renewcommand\tablename{Таблиця}}

% Додатки

\newcommand\myappendixes{
\section*{
    \begin{center}
    \thispagestyle{empty}
    \vspace*{\fill}
        \Large{Додатки}
        \addcontentsline{toc}{section}{\tocsecindent{Додатки}}
    \vspace*{\fill}
    \clearpage
    \end{center}
    }
\appendix
%\refstepcounter{section}

}

% Определение Додатків так, как нам надо ))
%\newcommand\myappendix[1]{ 
%    \refstepcounter{section} 
%	\section*{
%    \begin{center}
%      \thispagestyle{empty}
%      \vspace*{\fill}
%        Додаток~\thesection{}\\
%        #1
%    \addcontentsline{toc}{section}{Додаток~\thesection{}.~#1}
%      \vspace*{\fill}
%    \clearpage
%    \end{center}
%    } 
%}

\newcommand\myappendix[1]{ 
    \refstepcounter{section} 
	\section*{
    \begin{center}
        Додаток~\thesection{}. #1 \addcontentsline{toc}{section}{Додаток~\thesection{}.~#1}
    \end{center}
    } 
}

\makeatother
% Тело документа
\begin{document}

%ddcontentsline{toc}{section}{Приложение~\thesection{~#1}}}<++>Титульна сторінка.
\include{title}
%\maketitle % титулка гененрируемая классом. Переделывать пока лень, потому вставил вручную: см выше.
\newpage
\thispagestyle{empty}
\mbox{}
\newpage
\newpage
\thispagestyle{empty}
\mbox{}
\newpage
%
% Зміст
\renewcommand{\contentsname}{%
    \begin{center}\normalfont\normalsize Зміст\end{center}}
\tableofcontents
\newpage

% Розділи
\section*{Вступ}

\LaTeX (вимовляється «латех») — мова розмітки даних та пакет макросів \TeX для високоякісного оформлення документів, створений Леслі Лампортом (англ. Leslie Lamport). Вважається стандартом де-факто для підготовки математичних і технічних текстів для публікації в наукових виданнях.

На відміну від текстових процесорів, особливу увагу в \LaTeX приділено відокремленню змісту статті від оформлення. \LaTeX пропонує засоби для підготовки структурованих документів — документів, автор яких має можливість основну свою увагу зосередити на змісті, а оформлення і решту рутинної роботи перекласти на программу. Як і у випадку \TeX — вхідні файли \LaTeX можна порівняти із програмами.


\section{Розділ перший}
Цей документ своренно для тестування latex-стилю\cite{companion} для оформлення курсових робіт
за вимогами ВНТУ.

На данний момент стиль підтримує тільки курсові роботи (ДСТУ~3008-95). В
майбутньому (наступному триместрі) планується підтримка курсових проектів
(ГОСТ~2.105-95), і написання документації.
\subsection{Підрозділ перший}
Список того, що необхідно зробити, чи закінчити:
\begin{enumerate}
	\item підтримка бібліографії;
	\item підтримка ненумерованих розділів (аннотація, вступ, і т.п.);
	\item підтримка додатків;
	\item шаблони обов'язкових сторінок (титульний лист, обов'язкові додатки).
\end{enumerate}

\subsubsection{Підпідрозділ перший}

Складний список згідно госту:
\begin{enumerate}
	\item перший рядок першого рівня містить достатьно довгий текст, що повинен
		  перенестися на нову стрічку;
	\item другий рядок першого рівня містить:
	\begin{enumerate}
		\item перший підрядок другого рядка;
		\item другий підрядок другого рядка також містить достатьно довгий
			  текст, що повинен перенестися на нову стрічку;
		\item третій підрядок другого рядка:
	\end{enumerate}
	\item третій рядок третього рівня;
	\item 4;
	\item 5;
	\item 6;
	\item 7;
	\item 8;
	\item 9;
	\item 10;
	\item 11;
	\item 12;
	\item 13;
	\item 14;
	\item 15;
	\item 16;
	\item 17;
	\item 18;
	\item 19;
	\item 20;
	\item 21;
	\item 22;
	\item 23;
	\item 24;
	\item 25;
	\item 26;
	\item 27;
	\item 28;
	\item 29;
	\item 30;
 	\item 31.
% 	\item 32;
% 	\item 33;
% 	\item 34;
% 	\item 35;
% 	\item 36;
% 	\item 37;
%	\item 38;
%	\item 39;
%	\item 40.
\end{enumerate}

Текст після списку. Текст після списку. Текст після списку. Текст після списку. Текст після списку. Текст після списку. Текст після списку.

\begin{figure}[h]
  \centering
  \includegraphics{example}
  \caption{\label{img:img_example}Пример рисунка.}
\end{figure}

Метод Голда\cite{fstbook} получен из метода ван Циттерта\cite{2book}, но
с предположением, что действует он только для положительных чисел.
Восстанавливает по следующему алгоритму\cite{4book,5book,6book}.

Порядок методу дорівнює $p$, якщо існує таке позитивне число $c$, що
\begin{equation}\label{eq:delta}
  \Delta \le ch^{p+1},
\end{equation}

де ${\Delta}$~--- локальна похибка на кроці;

${h}$~--- крок дискретизації;

${p}$~--- порядок методу.

Текст після формули. Текст після формули. Текст після формули. Текст після формули. Текст після формули.

\section{Розділ перший}
Цей документ своренно для тестування latex-стилю\cite{companion} для оформлення курсових робіт
за вимогами ВНТУ.

На данний момент стиль підтримує тільки курсові роботи (ДСТУ~3008-95). В
майбутньому (наступному триместрі) планується підтримка курсових проектів
(ГОСТ~2.105-95), і написання документації.
\subsection{Підрозділ перший}
Список того, що необхідно зробити, чи закінчити:
\begin{enumerate}
	\item підтримка бібліографії;
	\item підтримка ненумерованих розділів (аннотація, вступ, і т.п.);
	\item підтримка додатків;
	\item шаблони обов'язкових сторінок (титульний лист, обов'язкові додатки).
\end{enumerate}

\subsubsection{Підпідрозділ перший}

Складний список згідно госту:
\begin{enumerate}
	\item перший рядок першого рівня містить достатьно довгий текст, що повинен
		  перенестися на нову стрічку;
	\item другий рядок першого рівня містить:
	\begin{enumerate}
		\item перший підрядок другого рядка;
		\item другий підрядок другого рядка також містить достатьно довгий
			  текст, що повинен перенестися на нову стрічку;
		\item третій підрядок другого рядка:
	\end{enumerate}
	\item третій рядок третього рівня;
	\item 4;
	\item 5;
	\item 6;
	\item 7;
	\item 8;
	\item 9;
	\item 10;
	\item 11;
	\item 12;
	\item 13;
	\item 14;
	\item 15;
	\item 16;
	\item 17;
	\item 18;
	\item 19;
	\item 20;
	\item 21;
	\item 22;
	\item 23;
	\item 24;
	\item 25;
	\item 26;
	\item 27;
	\item 28;
	\item 29;
	\item 30;
 	\item 31.
% 	\item 32;
% 	\item 33;
% 	\item 34;
% 	\item 35;
% 	\item 36;
% 	\item 37;
%	\item 38;
%	\item 39;
%	\item 40.
\end{enumerate}

Текст після списку. Текст після списку. Текст після списку. Текст після списку. Текст після списку. Текст після списку. Текст після списку.

\begin{figure}[h]
  \centering
  \includegraphics{example}
  \caption{\label{img:img_example}Пример рисунка.}
\end{figure}

Метод Голда\cite{fstbook} получен из метода ван Циттерта\cite{2book}, но
с предположением, что действует он только для положительных чисел.
Восстанавливает по следующему алгоритму\cite{4book,5book,6book}.

Порядок методу дорівнює $p$, якщо існує таке позитивне число $c$, що
\begin{equation}\label{eq:delta}
  \Delta \le ch^{p+1},
\end{equation}

де ${\Delta}$~--- локальна похибка на кроці;

${h}$~--- крок дискретизації;

${p}$~--- порядок методу.

Текст після формули. Текст після формули. Текст після формули. Текст після формули. Текст після формули.
\include{ch3}
\include{ch4}
\include{ch5}
\include{ch6}
%
\section{Розділ перший}
Цей документ своренно для тестування latex-стилю\cite{companion} для оформлення курсових робіт
за вимогами ВНТУ.

На данний момент стиль підтримує тільки курсові роботи (ДСТУ~3008-95). В
майбутньому (наступному триместрі) планується підтримка курсових проектів
(ГОСТ~2.105-95), і написання документації.
\subsection{Підрозділ перший}
Список того, що необхідно зробити, чи закінчити:
\begin{enumerate}
	\item підтримка бібліографії;
	\item підтримка ненумерованих розділів (аннотація, вступ, і т.п.);
	\item підтримка додатків;
	\item шаблони обов'язкових сторінок (титульний лист, обов'язкові додатки).
\end{enumerate}

\subsubsection{Підпідрозділ перший}

Складний список згідно госту:
\begin{enumerate}
	\item перший рядок першого рівня містить достатьно довгий текст, що повинен
		  перенестися на нову стрічку;
	\item другий рядок першого рівня містить:
	\begin{enumerate}
		\item перший підрядок другого рядка;
		\item другий підрядок другого рядка також містить достатьно довгий
			  текст, що повинен перенестися на нову стрічку;
		\item третій підрядок другого рядка:
	\end{enumerate}
	\item третій рядок третього рівня;
	\item 4;
	\item 5;
	\item 6;
	\item 7;
	\item 8;
	\item 9;
	\item 10;
	\item 11;
	\item 12;
	\item 13;
	\item 14;
	\item 15;
	\item 16;
	\item 17;
	\item 18;
	\item 19;
	\item 20;
	\item 21;
	\item 22;
	\item 23;
	\item 24;
	\item 25;
	\item 26;
	\item 27;
	\item 28;
	\item 29;
	\item 30;
 	\item 31.
% 	\item 32;
% 	\item 33;
% 	\item 34;
% 	\item 35;
% 	\item 36;
% 	\item 37;
%	\item 38;
%	\item 39;
%	\item 40.
\end{enumerate}

Текст після списку. Текст після списку. Текст після списку. Текст після списку. Текст після списку. Текст після списку. Текст після списку.

\begin{figure}[h]
  \centering
  \includegraphics{example}
  \caption{\label{img:img_example}Пример рисунка.}
\end{figure}

Метод Голда\cite{fstbook} получен из метода ван Циттерта\cite{2book}, но
с предположением, что действует он только для положительных чисел.
Восстанавливает по следующему алгоритму\cite{4book,5book,6book}.

Порядок методу дорівнює $p$, якщо існує таке позитивне число $c$, що
\begin{equation}\label{eq:delta}
  \Delta \le ch^{p+1},
\end{equation}

де ${\Delta}$~--- локальна похибка на кроці;

${h}$~--- крок дискретизації;

${p}$~--- порядок методу.

Текст після формули. Текст після формули. Текст після формули. Текст після формули. Текст після формули.

% Висновок
% TODO
\include{visnovok}
% Література
% TODO
\include{literatura}
% Додатки
\myappendixes

\myappendix{Алгоритм програми}
\input{images/prg-main-utf8.tpx}
\input{images/prg-e1-utf8.tpx}
\input{images/prg-e2-utf8.tpx}
\input{images/prg-print-utf8.tpx}



\include{listing}

\end{document}

